\documentclass[20pt,,margin=1in,innermargin=-4.5in,blockverticalspace=-0.25in]{tikzposter}
\geometry{paperwidth=43in,paperheight=32.5in}
\usepackage[utf8]{inputenc}
\usepackage{amsmath}
\usepackage{amsfonts}
\usepackage{amsthm}
\usepackage{amssymb}
\usepackage{mathrsfs}
\usepackage{graphicx}
\usepackage{adjustbox}
\usepackage{enumitem}
\usepackage{wrapfig}
\usepackage[backend=biber,style=numeric]{biblatex}
\usepackage{SUtheme}

\usepackage{mwe} % for placeholder images

\usepackage{comment}

\addbibresource{refs.bib}

% set theme parameters
\tikzposterlatexaffectionproofoff
\usetheme{SUTheme}
\usecolorstyle{SUStyle}
\usetitlestyle{Filled}

\usepackage[scaled]{helvet}
\renewcommand\familydefault{\sfdefault} 
\usepackage[T1]{fontenc}


\title{Lie Algebra of a Lie Group}
\author{Zih-Yu Hsieh, mentored by Arthur Jiang}
\institute{University of California Santa Barbara}
\titlegraphic{\includegraphics[width=0.06\textwidth]{logo.png}}

% begin document
\begin{document}
\maketitle
\centering
\begin{columns}
    \column{0.32}
    \block{Tangent Space, Tangent Vectors and Derivations}{
        %Delete everything about geometric tangent space (just a small intuition)
        In simplest case, if embedd manifold $M^n$ into $\mathbb{R}^m$, for any chart $(U,\phi)$ of $M$, since $\phi:U\rightarrow\phi(U)\subseteq \mathbb{R}^n$ has its inverse $\phi^{-1}$ being smooth, for any $u\in U\subseteq M$, a tangent vector $v_u$ associates with vector $v\in\mathbb{R}^n$, is characterized by differential of $\phi^{-1}$:
        $$v_u := D\phi^{-1}(\phi(u))(v) = \lim_{t\rightarrow 0}\frac{\phi^{-1}(\phi(u)+tv)-\phi^{-1}(\phi(u))}{t}$$
        A collection of all such vector is the \textbf{Geometric Tangent Space} of $u$, denoted as $T_u(M)$.

        \textbf{insert image}

        Notice that for any smooth function $f\in C^\infty(M)$, it has a notion of directional derivative at $u$ depending on the tangent vector $v_u\in T_u(M)$, and such derivative satisfies genral differentiation rules (for instance, product rule). 
        
        \hfil

        To generalize such notion into abstract manifold (space with no definition of vectors), we need a notion of \textbf{Derivation}: For any point $u\in M$, a \textbf{Derivation at $u$}, is a linear map $v_u:C^\infty(M)\rightarrow\mathbb{R}$, that satisfies the product rule:
        $$\forall f,g\in C^\infty(M),\quad v_u(fg) = f(u)(v_u g)+g(u)(v_u f)$$ 
        Which, the set of all derivations at $u$, denoted as $T_u(M)$, is the \textbf{Tangent Space} of $M$ at $u$, and each derivation $v_u\in T_u(M)$ is called the \textbf{Tangent Vector} of $u$.
    }
    \block{Vector Fields \& Smooth Conditions}{
        Given smooth manifold $M$, a vector field $X$ is a function associating each point $u\in M$ with a tangent vector of $u$, so $X(u)\in T_u(M)$. More precisely, a vector field is a map $X:M\rightarrow TM$ (where $TM$ denotes the \textbf{Tangent Bundle} of $M$), such that with the canonical projection map $\pi:TM\rightarrow M$, $\pi\circ X:M\rightarrow M$ is an identity.

        Which, $X$ is a \textbf{Smooth Vector Field}, if $X:M\rightarrow TM$ is a smooth map. And, a collection of smooth vector fields on $M$ is denoted as $\mathfrak{X}(M)$.
        
        \textbf{insert image}

        An equivalent condition of saying a vector field $X$ is smooth, is through smooth functions $f\in C^\infty(M)$: Since for all $u\in M$, $X(u)= X_u\in T_u(M)$ is a derivation at $u$, define $Xf:M\rightarrow\mathbb{R}$ by $Xf(u) = X_u(f)$. Then, $X$ is a smooth vector field iff $Xf\in C^\infty(M)$.
    }
    \block{Vector Fields of Different Manifolds}{
        Given $M,N$ two smooth manifolds, and smooth map $F:M\rightarrow N$. Let $X\in\mathfrak{X}(M)$, it would be ideal if we can send vector field $X$ to be a vector field of $N$. Yet, this requires both injectivity and surjectivity, which is too much to assume.
        
        \textbf{Insert an example}

        So, we'll consider a weaker notion, called an \textbf{$F$-Relation}: Given $X\in\mathfrak{X}(M)$ and $Y\in\mathfrak{X}(N)$, the two are $F$-related, if for all $u\in M$, the following is true:
        $$dF_u(X_u) = Y_{F(u)}$$
        Simply speaking, $F$ maps the tangent vectors collected by $X$, to be compatible with tangent vectors collected by $Y$.

        \textbf{Insert another example}

        \textbf{Thm:} If $F$ is a diffeomorphism, then for every $X\in\mathfrak{X}(M)$, there exists a unique $Y\in\mathfrak{X}(N)$, such that $X$ and $Y$ are $F$-related.
    }

    \column{0.36}

    \block{Lie Brackets on Vector Fields}{
        The initial motivation is to combine two vector fields $X,Y\in \mathfrak{X}(M)$ to be another vector field. Which, for all $f\in C^\infty(M)$, since $Yf\in C^\infty(M)$ from previous characterization, then $XYf = X(Yf)\in C^\infty(M)$. But, if consider function $XY$, in general it's not a vector field.

        EX: Define vector fields $X=\frac{\partial}{\partial x}$, $y=x\frac{\partial}{\partial y}$ on $\mathbb{R}^2$. Take smooth functions $f(x,y)=x$ and $g(x,y)=y$, then we get the following:
        $$XY(fg) = X(x\frac{\partial}{\partial y}(xy)) = \frac{\partial}{\partial x}(x^2) = 2x$$
        But, recall that vector field maps each point on $M$ to a derivation, so product rule should hold:
        $$XY(fg) = f(XY g) + g(XY f)$$
        With some computation, this equation doesn't hold for the example:
        $$f(XY g)+g(XY f)=x(X(x\frac{\partial}{\partial y}(y))) + y(X(x\frac{\partial}{\partial y}(x)))$$
        $$ = x(\frac{\partial}{\partial x}x)+y(X(x\cdot 0)) = x$$
        So, we need to define a new operation, called \textbf{Lie Bracket}, which is defined as:
        $$\forall X,Y\in\mathfrak{X}(M),\quad [X,Y]=XY-YX$$
        $$\forall f\in C^\infty(M),\quad [X,Y]f = X(Yf)-Y(Xf)$$
        Which, the output $[X,Y]\in\mathfrak{X}(M)$, and also satisfies the following:
        \begin{center}
            \begin{itemize}
                \item $\textbf{Bilinearity:}\quad [aX+bY,Z]=a[X,Z]+b[Y,Z]$
                \item $\textbf{Antisymmetry:}\quad [X,Y]=-[Y,X]$
                \item $\textbf{Jacobi's Identity:}\quad \left[X,[Y,Z]\right]+ \left[Y,[Z,X]\right]+ \left[Z,[X,Y]\right]=0$
            \end{itemize}
        \end{center}
        Moreover, given smooth map $F:M\rightarrow N$, if $X_1,X_2\in\mathfrak{X}(M)$ and $Y_1,Y_2\in\mathfrak{X}(N)$ are $F$-related respectively, then $[X_1,X_2]\in \mathfrak{X}(M)$ and $[Y_1,Y_2]\in\mathfrak{X}(N)$ are also $F$-related. This is the essential tool for defining Lie Algebra on a Lie Group.
    }
    \block{Lie Group \& Left-Invariant Vector Fields}{
        A \textbf{Lie Group} $G$, is a smooth manifold along with group structure, such that the group operation $P:G\times G\rightarrow G$ by $P(g,h) = gh$, and the inversion map $i:G\rightarrow G$ by $i(g)=g^{-1}$ are both smooth maps between manifolds.

        \textbf{EX: eucildean space}

        For all $g\in G$, denote the left multiplication $L_g:G\rightarrow G$ by $L_g(h)=gh$,
        since $L_g = P\bigm|_{\{g\}\times G}$, all left multiplication is a smooth map; also, since $L_{g^{-1}}\circ L_g(h) = L_{g^{-1}}(gh) = g^{-1}gh = h$, every left multiplication has a smooth inverse, hence it's a \textbf{Diffeomorphism}.

        \hfil

        \textbf{Left-Invariant vector fields:} Given any $X\in\mathfrak{X}(G)$ and all $g\in G$, since $L_g$ is a diffeomorphism, there's a notion of $X$ being $L_g$-related to itself. Which, we call $X$ a \textbf{Left-Invariant Vector Field}, if for all $g\in G$, $X$ is $L_g$-related to itself.

        \textbf{EX: euclidean addition, circle or torus}

        \hfil

        Recall that Lie Bracket of vector field preserves an $F$-relation. Then, for all $X,Y\in\mathfrak{G}$ that are left-invariant, since for all $g\in G$, $X$ and $Y$ are $L_g$ related to themselves, then the Lie Bracket $[X,Y]$ is also $L_g$ related to $[X,Y]$. Hence, $[X,Y]$ is also left-invariant, Lie Bracket on a Lie Group preserves left invariance of vector fields.
    }

    \column{0.32}
    \block{Lie Algebra on a Lie Group}{
        
    }

    \block{Application}{

    }
    
   \block{Acknowledgements}{
        I want to thank my mentor Arthur Jiang for the effort and kindness of guiding me through the materials, and provide helpful information on the project. I'd also like to thank the UCSB Math Department Directed Reading Program for this opportunity.
    }
    
    \block{References}{
        John M. Lee, Introduction to Smooth Manifolds, 2nd Edition
    }
\end{columns}
\end{document}